\documentclass{article}
\usepackage{amsmath, amssymb, amsthm}
\usepackage{hyperref}
\usepackage{geometry}
\geometry{margin=1in}

\title{Erd\H{o}s Problem \#593}
\date{Last updated: 2025-08-31}
\author{Source: erdosproblems.com}

\begin{document}
\maketitle

\noindent\textbf{Status:} OPEN \quad \textbf{Prize: \$500}

\noindent\textbf{Tags:} set theory, graph theory, hypergraphs, chromatic number

\medskip
\noindent\textbf{Problem Statement:}

Characterize those finite 3-uniform hypergraphs which appear in every 3-uniform hypergraph of chromatic number $>\aleph_0$.

\bigskip
\noindent\textbf{Remarks:}

Similar problems were investigated by Erd\H{o}s, Galvin, and Hajnal \cite{EGH75}. Erd\H{o}s claims that for graphs the problem is completely solved: a graph of chromatic number $\geq \aleph_1$ must contain all finite bipartite graphs but need not contain any fixed odd cycle.

\bigskip
\noindent\textbf{References:}

[EGH75] Erd\H{o}s, P. and Galvin, F. and Hajnal, A., On set-systems having large chromatic number and not containing prescribed subsystems. (1975), 425--513.

\bigskip
\noindent\small{Source: \url{https://www.erdosproblems.com/593}}
\end{document}
